\chapter{拓扑结构}

\section{开集}

记号约定:$X,Y$是集合.

\begin{definition}{拓扑,开集}{}
	设$\mathfrak{O}\subseteq\mathcal{P}\left(X\right)$.若$\mathfrak{O}$满足下列条件,则称之为$X$上的拓扑,$\mathfrak{O}$的元素称为$X$中(相对$\mathfrak{O}$)的开集.
	\begin{enumerate}
		\item ($\mathrm{O}_{1}$)$\forall\mathfrak{S}\subseteq\mathfrak{O}\left(\bigcup\limits_{S\in\mathfrak{S}}S\in\mathfrak{O}\right)$.
		\item ($\mathrm{O}_{2}$)$\forall\mathfrak{S}\in\fin\left(\mathfrak{O}\right)\left(\bigcap\limits_{S\in\mathfrak{S}}S\in\mathfrak{O}\right)$.
	\end{enumerate}
\end{definition}

\begin{definition}{拓扑空间}{}
	配备了拓扑的$X$称为拓扑空间,$X$中的元素通常被称为点.
\end{definition}

\begin{remark}{验证拓扑的办法}{}
	根据上述定义,$\emptyset$和$X$都属于$\mathfrak{O}$.验证$\mathfrak{O}$是否满足($\mathrm{O}_{2}$)等价于验证下列条件同时成立:
	\begin{enumerate}
		\item ($\mathrm{O}_{2a}$)$\forall S_{1},S_{2}\in\mathfrak{O}\left(S_{1}\cap S_{2}\in\mathfrak{O}\right)$.
		\item ($\mathrm{O}_{2b}$)$X\in\mathfrak{O}$.
	\end{enumerate}
\end{remark}

\begin{example}{拓扑的例子}{}
	\begin{enumerate}
		\item $\left\{\emptyset,X\right\}$是$X$上的拓扑,称为平庸拓扑,配备离散拓扑的$X$称为平庸空间.
		\item $\mathcal{P}\left(X\right)$是$X$上的拓扑,称为离散拓扑,配备离散拓扑的$X$称为离散空间.
	\end{enumerate}
\end{example}

本节接下来的内容中,默认$X$和$Y$都配备了拓扑.

\begin{definition}{开覆盖}{}
	设$A\subseteq X$,而$\left(U_{i}\right)_{i\in I}$是$X$中的一族开集.若该开集族是$A$的覆盖,则称之为$A$在$X$中的开覆盖.
\end{definition}

\begin{definition}{同胚}{}
	设$f:X\to Y$是双射.若$X$中每个开集在$f$下的像都是$Y$中的开集,则称$f$是同胚.记$\Hom_{\mathsf{Top}}\left(X,Y\right)$是$X$到$Y$的同胚组成的集合.
\end{definition}

\begin{definition}{同胚的拓扑空间}{}
	若$\Hom_{\mathsf{Top}}\left(X,Y\right)\neq\emptyset$,则称$X$和$Y$同胚.
\end{definition}

\begin{theorem}{同胚的范畴版定义}{}
	设$f:X\to Y$是双射,则$f$是同胚当且仅当下列条件成立:
	\begin{enumerate}
		\item $X$中每个开集在$f$下的像都是$Y$中的开集.
		\item $Y$中每个开集在$f^{-1}$下的像都是$X$中的开集.
	\end{enumerate}
\end{theorem}


























